\subsubsection{Herstellung des Gehäuses}
Zur Herstellung des Gehäuses kommt ein 3D-Drucker der Marke Creality 3D zum Einsatz. Der stark modifizierte Ender 3 Pro von Manuel Starz (vgl. \ref{fig:ender3}) druckt die zuvor erstellten STL-Dateien mit PETG-Filament der Firma dasfilament. Der für den 3D-Druck nötige G-Code wird mit Ultimaker CURA generiert und mit Hilfe von OctoPrint (vgl. \ref{fig:octoprint}) an den Drucker übertragen. Das Gehäuse besteht aus vier Einzelteilen, je zwei Teile für die Seitenwände und zwei Teile für den Deckel.\par
\begin{figure}[h!tb]
	\begin{center}
	\includegraphics[width=.5\textwidth,angle=270]{img/ender3pro_mod.jpg}
	\end{center}
	\caption[3D-Drucker Ender 3 Pro (mod)]{3D-Drucker Ender 3 Pro (mod)}
	\label{fig:ender3}
\end{figure}
\begin{figure}[h!tb]
	\includegraphics[width=1\textwidth]{img/gehäusedruck_octoprint_01.png}
	\caption[OctoPrint-Weboberfläche]{OctoPrint-Weboberfläche}
	\label{fig:octoprint}
\end{figure}
Die gedruckten Teile (vgl. \ref{fig:printet_parts}) werden dann zum Teil mit Zwei-Komponenten-Epoxidkleber verbunden (vgl. \ref{fig:glued_parts_01} \& \ref{fig:glued_parts_02}) und mit Hilfe von Schleifpapier (vgl. \ref{fig_filed_parts}) und einigen Schichten Klarlack zu einem klavierlackähnlichen Finish veredelt (vgl. \ref{fig:finished_parts}). Die Seitenwände werden dann mit dem Bildschirm mit Hilfe des Zwei-Komponenten-Epoxidklebers permanent verklebt. Die Rückseite besteht der Einfachheit halber aus zwei identischen, punktsymmetrischen Teilen, die ebenfalls miteinander verklebt wurden.\par
\begin{figure}[h!tb]
	\includegraphics[width=1\textwidth]{img/placeholder.png}
	\caption[Ausgedruckte Teile des Gehäuses]{Ausgedruckte Teile des Gehäuses}
	\label{fig:printet_parts}
\end{figure}
\begin{figure}[h!tb]
	\includegraphics[width=1\textwidth]{img/placeholder.png}
	\caption[Verklebung des Gehäusedeckels]{Verklebung des Gehäusedeckels}
	\label{fig:glued_parts_01}
\end{figure}
\begin{figure}[h!tb]
	\includegraphics[width=1\textwidth]{img/placeholder.png}
	\caption[Verklebung der Gehäuseseiten]{Verklebung der Gehäuseseiten}
	\label{fig:glued_parts_02}
\end{figure}
\begin{figure}[h!tb]
	\includegraphics[width=1\textwidth]{img/placeholder.png}
	\caption[Abschleifen der Gehäuseaußenseite]{Abschleifen der Gehäuseaußenseite}
	\label{fig:filed_parts}
\end{figure}
\begin{figure}[h!tb]
	\includegraphics[width=1\textwidth]{img/placeholder.png}
	\caption[Fertig lackiertes Gehäuse]{Fertig lackiertes Gehäuse}
	\label{fig:finished_parts}
\end{figure}
Laut Ultimaker CURA beträgt die Gesamtdruckdauer des Gehäuses 36 Stunden 36 Minuten und verbraucht insgesamt 423 Gramm des verwendeten Filaments (dasFilament PETG Schwarz). Damit belaufen sich die Materialkosten des Gehäuses auf 11,61\euro{}.
\newpage